% TOP_PROBLEM: {"id": 1509, "title": "Siegel Zeros", "category_id": 1, "category": {"id": 1, "name": "number_theory", "display_name": "Number Theory", "description": "Properties of integers, prime numbers, Diophantine equations.", "slug": "number-theory", "order_index": 1, "created_at": "2026-07-31T15:26:25.670Z"}, "status": "open"}
% FIRST_POSTED: 2026-10-19
% ATTEMPT_STATUS: unresolved
% SUBMISSION: {"kind": "research_attempt", "category": "top_problems", "problem_id": "1509", "model": {"provider": "OpenAI", "version": "GPT 6 Astra Pro (exact release/build unknown)", "reasoning": "Unknown; not inferred from product name", "generated_on": "2026-10-19"}, "other_models": "None separately identified for the supplied report; original model build is unknown. Repository formatting performed by Codex (exact underlying build not recorded).", "tools": "Source consultation and supplied images as reported in manuscript; Pandoc for Markdown-to-LaTeX integration. No theorem prover or new computation.", "target": "Uniform c>0 for every primitive real nonprincipal character of conductor q>=3: every real zero beta in (0,1) satisfies (1-beta) log q>=c.", "scope": "Comparative literature exposition only; no new theorem, independent certification or status change.", "human_contributions": "Submitter supplied manuscript images and provenance URL and requested comparison; no human mathematical authorship or independent checking represented.", "verification": "Source checks reported by original submission; integration checks formatting and scope consistency only, not proofs or image authenticity.", "reproducibility": "Report source-access date 2026-10-19 used as document date; exact generation date unknown. Original prompts and images not included; initial/final research-prompt use not documented. Markdown converted with Pandoc; literature-report caveats preserved."}
\documentclass[11pt]{article}
\usepackage[T1]{fontenc}
\usepackage[utf8]{inputenc}
\usepackage[margin=25mm]{geometry}
\usepackage{amsmath,amssymb}
\usepackage{hyperref}
\usepackage{bookmark}
\hypersetup{hidelinks}
\setlength{\emergencystretch}{3em}
\newcommand{\tightlist}{\setlength{\itemsep}{0pt}\setlength{\parskip}{0pt}}
\title{From interpolation determinants to Wronskians: a literature report on Siegel zeros}
\author{OpenAI GPT 6 Astra Pro}
\date{19 October 2026}
\begin{document}
\maketitle
\section{Scope and identity}\label{scope-and-identity}

Canonical problem: 1509, Siegel Zeros (We-Solve Math-77). Claim type: literature report.

Comparative exposition of the OpenAI, Banks, Kerr, Green and Barreto manuscripts on the uniform exclusion of Landau--Siegel zeros, with particular attention to the three-page manuscript bearing Kevin Barreto's byline. The report explains the Wronskian nonvanishing argument, its interaction with weighted monomials and prime divisibility, and the subsequent passage from a prime-sum estimate to the real-zero bound. It supplements \href{https://we-solve.github.io/contributions/siegel-zeros/attempt-2/r1/}{Math-77-A2} with a comparison of subsequent methods and a detailed account of Barreto's simplification.

This is an exposition of the cited arguments, not a claim of a new theorem or an independent certification of the five proofs. Public sources were consulted for OpenAI, Banks, Kerr and Green; Barreto's mathematical text was read from the three supplied page images. The submitter identifies \href{https://x.com/AcerFur/status/2112239827077026128}{this post by @AcerFur} as the document's sharing source. The post itself could not be retrieved during this review; its URL is recorded as supplied provenance, not as independent authentication of the images. No change to Math-77's verification status is requested.

\subsection{1. The common problem and the purpose of this report}\label{the-common-problem-and-the-purpose-of-this-report}

The target in \href{https://we-solve.github.io/problems/siegel-zeros/}{We-Solve's Math-77} is a \textbf{single constant that works for every conductor}. In the notation used here, it is

\[
\begin{gathered}
\exists c>0\quad \forall q\ge3\quad
\forall\chi\text{ primitive, real and nonprincipal of conductor }q:\\
L(\beta,\chi)=0,\quad 0<\beta<1
\quad\Longrightarrow\quad (1-\beta)\log q\ge c.
\end{gathered}
\]

All five manuscripts concern this assertion. It excludes exceptionally close real zeros near one; it does not exclude every real zero in the critical strip or establish the Generalized Riemann Hypothesis. A separate positive gap for each character would not suffice.

The existing \href{https://we-solve.github.io/contributions/siegel-zeros/attempt-2/r1/}{Math-77-A2 literature report} documents OpenAI's claims. This contribution concentrates on how the subsequent presentations change the algebraic argument, especially the replacement of Green's resultant calculation by Barreto's Wronskian. References below distinguish sources read directly from background results consulted through those sources.

\subsection{2. What changes across the five presentations}\label{what-changes-across-the-five-presentations}

\textbf{OpenAI: a weighted interpolation determinant.} The manuscript dated 1 October 2026 {[}OAI{]} works, apart from a fixed exceptional field handled separately, in \(\mathbb Q(\sqrt d,\sqrt2)\). A four-dimensional integer box supplies points in three conjugate coordinates. An interpolation lemma establishes enough rank, and a greedy weighted selection produces a nonzero square determinant. Frobenius congruences at primes with \(\chi(p)=-1\) give divisibility incompatible with its size when a hypothetical zero forces sufficient prime bias. The key scale comparison is between the number of points and the degree required for interpolation.

The determinant's entries are \textbf{monomial evaluations}; the derivative matrices below should not be attributed to this construction. The paper also discusses an interpretation in terms of derivatives, but that is distinct from saying its displayed evaluation matrix has derivative entries. Its absolute zero-gap constant is not given numerically. See \href{https://github.com/openai/math/blob/adc7f1241b42e322a6451854ab7e4b4c146bf78a/preprints/Uniform-exclusion-of-Landau-Siegel-zeros-October-1-2026/build/paper.tex}{OAI, Sections 2--6}.

\textbf{Banks: explicit constants and an expanded direct proof.} Banks {[}Ba{]} retains the biquadratic interpolation method but reorganizes the proof and tracks its estimates quantitatively. His Theorem 1.1 states the admissible values \(q_0=10^{10}\) and \(\delta_0=0.003\). For the smaller conductors, he invokes the computation of Lu--Zaman--Zhao {[}LZZ{]}, whose Theorem 1.1 excludes real zeros in \(\sigma\ge1-1/(5\log q)\) for \(q\le10^{10}\). The resulting stated bound is \((1-\beta)\log q\ge0.003\) for all \(q\ge3\).

The contribution is principally explicitness and exposition, not a replacement of the algebraic mechanism. Banks credits OpenAI with the underlying ideas, describes Claude's assistance, and states that he personally checked the entire argument. That is the author's verification statement, not a check independently performed by this report. See \href{https://arxiv.org/html/2610.15483v1}{Ba, introduction and Theorem 1.1}.

\textbf{Kerr: replace determinant analysis by pigeonholing.} Kerr {[}Ke{]} keeps the biquadratic setting but counts bounded-height polynomials instead of selecting an interpolation determinant. A pigeonhole argument produces a nonzero polynomial subject to strong modular constraints. An algebraic-norm bound then upgrades divisibility to exact vanishing on the projected integer box. A degree lower bound attributed to Masser {[}Ma{]}, stated as Kerr's Lemma 6, contradicts the degree of this polynomial.

Thus the replacement is not simply a shorter determinant computation: the auxiliary object becomes a vanishing polynomial, and the final obstruction is a polynomial nonvanishing theorem. Kerr credits ChatGPT assistance and Banks's exposition. See \href{https://arxiv.org/html/2610.21099v1}{Ke, Lemmas 5--7 and the concluding proof}. Masser's original theorem is not reproved or independently audited here.

\textbf{Green: quadratic-field jets and resultants.} Green {[}Gr{]} removes the auxiliary \(\sqrt2\). Derivatives in one variable of bivariate monomials, evaluated at quadratic-conjugate pairs, form rows of an integral Gram-matrix construction. Row independence is proved by polynomial zero counting and a resultant bound for specialized root multiplicities; determinant size and prime divisibility then give the prime-sum estimate. Version 2 separates the lattice-size and derivative-order parameters and states, for \(X=q^C\),

\[
\sum_{\substack{p\le X\\\chi(p)=-1}}\frac{\log p}{p}
\le\left(\frac12+O(1/C)\right)\log X+O_C(1).
\]

Importantly, \textbf{Green's version 2 already remarks that a nonsingular minor can replace the Gram determinant}. That alternative must not be presented as a new idea originating with Barreto. Green's provenance statement credits ChatGPT with the arguments and himself with the exposition. See \href{https://hexagonmath.org/pdf/2610.00311v2}{Gr, Theorem 1.1, the remark before Section 3, and Appendix A}.

\textbf{Barreto: the Wronskian replacement.} The supplied three-page manuscript {[}Br{]} expressly replaces Green's resultant-based independence argument by an ordinary polynomial Wronskian. It combines this with a weighted triangular set of monomials and a rectangle chosen to keep both quadratic conjugates small. The next section explains why these choices work together, rather than merely observing that a Wronskian appears in the proof.

\subsection{3. Barreto's three-page argument: the essential steps}\label{barretos-three-page-argument-the-essential-steps}

The manuscript has two parts: \textbf{Theorem 1 bounds a prime sum}, and \textbf{Corollary 2 converts that bound into a restriction on real zeros}. The first part does not assume that an exceptional zero exists. This separates the algebraic estimate from its analytic application. All implied constants in the manuscript are absolute.

Its Theorem 1 states, uniformly for \(x\ge2\),

\[
A_\chi(x):=\sum_{\substack{p\le x\\\chi(p)=-1}}\frac{\log p}{p}
\le \frac12\log x+\frac14\log q+\log\log(3x)+O(1).
\tag{1}
\]

The nontrivial point is the coefficient \(1/2\) of \(\log x\), with controlled conductor dependence. Bounding by the sum over all primes would instead give \(\log x+O(1)\). The account below follows {[}Br, pp.~1--3{]}, expanding the logical connections while omitting its routine summation calculations.

\textbf{3.1. Small conjugates and weighted monomials.} Associate \(\chi\) with a squarefree \(d\ne1\), \(|d|\le q\), so that \(\chi(p)=(d/p)\) for odd \(p\nmid q\). The manuscript chooses a rectangle of integer pairs \((u,v)\) of area \(2x\), balanced between the sizes of \(u\) and \(v\sqrt q\). This supplies at least \(2x\) points

\[
\alpha_i=u+v\sqrt d,\qquad \bar\alpha_i=u-v\sqrt d,
\qquad |\alpha_i|,|\bar\alpha_i|<B,
\qquad B=2\sqrt{2x}\,q^{1/4}.
\]

Here the bar means \textbf{quadratic-field conjugation}, not necessarily complex conjugation. Each of the two lists is pairwise distinct because \(\sqrt d\) is irrational. The balancing matters quantitatively: \(\log B=\tfrac12\log x+\tfrac14\log q+O(1)\), which will supply the principal terms in (1).

Give \(z\) weight \(1\) and \(w\) weight \(k\), where \(3\le k\le x\) is chosen later. Use the monomials indexed by

\[
T=\{(r,s)\in\mathbb Z_{\ge0}^2:r+ks<x\}.
\]

The row for \((r,s)\) consists of \(\partial_z^j(z^rw^s)(\alpha_i,\bar\alpha_i)\) for every point and every integer \(0\le j<2x/k\). These are ordinary derivatives, sometimes called jets: the value and its successive derivatives. Splitting each entry into its two integer coefficients in the basis \(1,\sqrt d\) gives an integer matrix.

There are two reasons for this setup. A small degree in \(w\) will control the dimension of a family of specialized polynomials. The weighted support will also remain closed under the substitutions needed for prime divisibility. The same parameter \(k\) therefore connects the rank argument to the arithmetic argument.

\textbf{3.2. Why a Wronskian proves row independence.} Suppose the integer rows were dependent. Since the matrix is rational, a dependence can be chosen over \(\mathbb Q\), producing a nonzero polynomial \(F(z,w)\) supported on \(T\). It has \(\deg_zF<x\) and \(\deg_wF<x/k\), and every specialization \(F(z,\bar\alpha_i)\) has all the prescribed derivatives zero at \(z=\alpha_i\).

Some specializations could vanish identically, but there are fewer than \(x/k\) such values of \(w\): any nonzero coefficient polynomial of \(F\), viewed as a polynomial in \(z\), already bounds their number. Consequently, more than \(x\) specializations are nonzero. Their span is a space \(V\) of polynomials in \(z\) of dimension \(\ell<x/k+1\).

Choose a monic basis \(g_1,\ldots,g_\ell\) of \(V\) with distinct degrees, by successively cancelling leading terms. Consider

\[
W(z)=\det\bigl(g_b^{(a)}(z)\bigr)_{
0\le a<\ell,\;1\le b\le\ell}.
\qquad W\ne0,\qquad \deg W<\ell x.
\tag{2}
\]

Why is it nonzero? For distinct degrees \(n_b=\deg g_b\), its leading-coefficient determinant has entries given by the falling factorials \(n_b(n_b-1)\cdots(n_b-a+1)\). These are monic polynomials of the successive degrees \(a=0,\ldots,\ell-1\), evaluated at distinct numbers \(n_b\). Their determinant is a nonzero Vandermonde determinant. This is a polynomial argument in characteristic zero; it does not invoke an unrestricted claim about Wronskians of arbitrary differentiable functions.

Now fix any nonzero specialization. Write \(m=\lceil2x/k\rceil\) for the number of derivative conditions. The specialization has a zero of order at least \(m\) at its associated \(\alpha_i\). Make it a member of a new basis of \(V\). Such a basis change multiplies the Wronskian by a nonzero constant. In the column corresponding to this basis member, derivatives through order \(\ell-1\) still contain the factor \((z-\alpha_i)^{m-\ell+1}\).

The dimension bound gives \(2\ell-2<2x/k\), hence \(m-\ell+1\ge\ell\). Therefore \textbf{the same Wronskian has a zero of multiplicity at least \(\ell\) at each of more than \(x\) distinct points}. Its degree would exceed \(\ell x\), contradicting (2).

This is the decisive simplification: instead of tracking multiplicities of irreducible factors through resultants, one packages the specializations into one finite-dimensional polynomial space and counts zeros of one nonzero Wronskian. The integer matrix consequently has full row rank, so a nonzero square integer minor \(D\) using all rows exists. No particular column-selection algorithm is needed.

\textbf{3.3. Why prime divisibility survives differentiation.} Let \(k\le p\le x\) and \(\chi(p)=-1\). These primes are odd and do not divide \(q\). Euler's criterion and the binomial theorem give \(\alpha_i^p\equiv\bar\alpha_i\pmod{p\mathbb Z[\sqrt d]}\). Thus \(\phi(z,w)=z^p-w\) takes values divisible by \(p\) at the chosen points.

The essential additional fact is that \textbf{every positive ordinary \(z\)-derivative of \(\phi\) also has coefficients divisible by \(p\)}; derivatives beyond its degree are zero. Consequently, differentiating a product of \(h\) copies of \(\phi\) does not destroy its \(p^h\) divisibility at these points. This explains why the derivative matrix remains compatible with Frobenius. It is not merely a congruence for the undifferentiated entries.

For \(h=\lfloor r/p\rfloor\), replace the row arising from \(z^rw^s\) by the row arising from

\[
z^{r-ph}(z^p-w)^h w^s
=z^rw^s+\sum_{j=1}^{h}(-1)^j\binom hj z^{r-pj}w^{s+j}.
\tag{3}
\]

Every added monomial stays in \(T\), since \(r-pj+k(s+j)=r+ks-(p-k)j\le r+ks<x\). It also has smaller \(z\)-degree. Ordering rows by that degree makes (3) an integer triangular row transformation with diagonal entries one. The chosen minor is unchanged, while its transformed \((r,s)\) row is divisible by \(p^{\lfloor r/p\rfloor}\). Therefore

\[
p^{\sum_{(r,s)\in T}\lfloor r/p\rfloor}\mid D.
\tag{4}
\]

The role of the weights is now visible: replacing a block \(z^p\) by \(w\) does not increase weight when \(p\ge k\). This is the support-preservation fact that lets the full exponent \(\lfloor r/p\rfloor\) be used.

\textbf{3.4. Comparing divisibility with size.} Put \(M=|T|\) and \(U=\sum_{(r,s)\in T}r\). The manuscript's lattice-point estimates give \(M\ll x^2/k\) and \(U\gg xM\). The latter relation absorbs two kinds of error.

First, replacing \(\lfloor r/p\rfloor\) by \(r/p\) loses at most \(M\) per prime. Chebyshev's bound makes the total logarithmic loss \(O(xM)=O(U)\). Second, determinant expansion and the derivative bounds also fit within \(O(U)\) after choosing the cutoff appropriately. Primes below \(k\) contribute at most \(\log k+O(1)\) by Mertens' estimate.

For the size estimate, the row indexed by \((r,s)\) is bounded using \(B^r(x^2B)^{x/k}\). Choosing \(k=\lceil\log(x^2B)\rceil\) limits the extra factor to \(e^x\) per row. The factorial in determinant expansion is likewise absorbed. The manuscript consequently obtains \(\log|D|\le U\log B+O(U)\). Combining this with (4) gives

\[
A_\chi(x)
\le\frac{\log|D|}{U}+\log k+O(1)
\le\log B+\log k+O(1).
\tag{5}
\]

For this step the manuscript first reduces to \(x\ge\max(e^3,\sqrt q)\), where \(3\le k\le x\) and \(\log k=\log\log x+O(1)\). The omitted small range follows directly from Mertens' estimate. Substitution of the value of \(B\) in (5) proves the stated bound (1).

Thus the conductor term comes from controlling both conjugates, the logarithmic-logarithmic term from the prime cutoff, and the saving in the main prime sum from comparing the divisibility of a nonzero integer with its size. These are explanations of the manuscript's choices, not additional claimed improvements.

\textbf{3.5. From the prime sum to real zeros.} Corollary 2 applies classical analytic number theory to (1). For \(0<\eta\le1\), partial summation gives

\[
S_\eta:=\sum_{\chi(p)=-1}\frac{\log p}{p^{1+\eta}}
\le\frac{1}{2\eta}+\frac14\log q+O(\log(2/\eta)).
\tag{6}
\]

The manuscript controls the integrated \(\log\log(3t)\) error by Jensen's inequality. Higher prime powers have a uniformly bounded contribution, while primes dividing \(q\) contribute \(O(\log\log(3q))\). The Euler products then give

\[
\frac{L'}{L}(1+\eta,\chi)
=2S_\eta-\frac1\eta+O(\log\log(3q))
\le\frac12\log q+O\bigl(\log(2/\eta)+\log\log(3q)\bigr).
\tag{7}
\]

The important cancellation is \(2\cdot(1/(2\eta))-1/\eta=0\). This explains the significance of the coefficient \(1/2\) in the prime-sum bound, rather than treating the passage to zeros as a black box.

For a real zero \(\beta\in(0,1)\), retain its contribution in the Hadamard-product formula. At the real point \(1+\eta>1\), the other nontrivial zeros contribute nonnegative real parts; the gamma term is uniformly bounded for both parities. Hence

\[
\frac{1}{1-\beta+\eta}
\le\frac{L'}{L}(1+\eta,\chi)+\frac12\log q+O(1)
\le\log q+O\bigl(\log(2/\eta)+\log\log(3q)\bigr).
\tag{8}
\]

Taking \(\eta=(\log q)^{-2}\), inverting the upper bound for sufficiently large \(q\), and subtracting \(\eta\) yields the manuscript's conclusion

\[
(1-\beta)\log q
\ge1-O\!\left(\frac{\log\log(3q)}{\log q}\right).
\tag{9}
\]

This gives a positive uniform gap for large conductors. There are only finitely many remaining characters, each with finitely many real zeros in the relevant compact interval, so decreasing the constant covers them as well. This final finiteness argument is not a new finite-range numerical computation.

The asymptotic leading constant in (9) must not be confused with a demonstrated numerical bound \(c=1\) for every conductor. Banks's all-conductor explicit constant and Barreto's asymptotic lower bound are different kinds of quantitative statement.

\subsection{4. What is actually simplified, and what is retained}\label{what-is-actually-simplified-and-what-is-retained}

The central replacement is \textbf{resultant-based multiplicity control by a Wronskian degree argument}. Barreto does not discard the determinant method: the Wronskian proves that an arithmetic determinant is nonzero, after which the same size-versus-divisibility principle remains essential.

There are useful accompanying design changes. Barreto chooses the minor route already mentioned by Green, uses coefficient splitting to work with an ordinary integer determinant, balances the rectangle of algebraic integers, and uses the weighted triangular support \(r+ks<x\). The algebra in (3) explains directly why that support is convenient. It would be misleading to credit every one of these ingredients as a separately new theorem; the report identifies how they are assembled in the supplied manuscript.

The contrast with Kerr is also precise. Kerr avoids the interpolation determinant by constructing an impossible vanishing polynomial and invoking Masser. Barreto retains a nonzero minor and proves the needed rank statement using elementary polynomial linear algebra. Neither description should be collapsed into the vague statement that both papers ``use a simpler determinant.''

Finally, ``elementary'' describes the streamlined algebraic portion and presentation, not the removal of the classical theory of Dirichlet \(L\)-functions. Barreto still uses Chebyshev and Mertens estimates, the quadratic-character correspondence, Euler products and the Hadamard-product formula, with references to Montgomery--Vaughan {[}MV{]}.

\subsection{5. Attribution, source checks and verification limits}\label{attribution-source-checks-and-verification-limits}

The supplied manuscript's opening paragraph credits GPT-6 Astra with the Wronskian simplification and assistance checking its text. This is the manuscript's own attribution. Likewise, the AI-assistance statements in Banks, Kerr and Green describe the preparation of those sources; they do not mean that those authors or models participated in preparing this literature report.

This revision was checked against Barreto's three page images, the pinned OpenAI TeX source, the Banks and Kerr HTML texts, Green's version-2 record and PDF, and the relevant We-Solve pages. The comparison includes local algebraic explanations of the Wronskian and triangular substitution, but it is not an independent expert audit of all estimates and background theorems. No Lean development was executed and no finite-conductor computation was reproduced. The reported formalization scope already discussed in \href{https://we-solve.github.io/contributions/siegel-zeros/attempt-2/r1/}{Math-77-A2} is not promoted here into an independently verified claim.

The contribution is therefore a source-based explanation of the progression of methods, with the Wronskian proof made intelligible and correctly attributed. The separate fixed-half-plane claims also listed in Math-77-A2 are outside its scope. No peer-review outcome, formal verification of Barreto's particular argument, or change in the site's status is asserted.

\section{Model attribution and preparation}\label{model-attribution-and-preparation}

The supplied report attributes its mathematical exposition to OpenAI GPT 6 Astra Pro. Exact release/build, reasoning setting and thinking time are unknown. The report is dated 19 October 2026 by its source-access statement; a separate generation timestamp was not supplied. The submitter supplied the manuscript images and provenance URL and requested the literature comparison; no human mathematical authorship or independent proof audit is asserted. No additional model is identified for the report. The source authors' model disclosures concern their manuscripts, not authorship of this report.

\section{Reproducibility and verification}\label{reproducibility-and-verification}

The source-access statements in this report describe the preparation of the supplied submission. Repository integration converts that submission from Markdown to LaTeX with Pandoc and adds identity, attribution and status metadata. The original images, browsing transcripts and original generation prompts are not distributed here. No fresh source authentication, independent mathematical proof verification, Lean execution or finite-conductor computation was performed during integration. The original use of the repository's initial and final research prompts is not documented. The integration review preserves the uniform all-conductor target and distinguishes exposition from a new result.

\section{Final status}\label{final-status}

UNRESOLVED as a new research attempt: this is a literature report, not a new theorem or independent certification. The cited authors' claims remain attributed to those sources; no source-catalogue or community-review status change is requested.

\section{Completion Estimate}\label{completion-estimate}

\noindent\textbf{COMPLETION: 0\%}

No new progress toward proving the original all-conductor assertion is claimed by this expository report. This estimate measures the report's original research contribution, not the extent of the results claimed in the cited literature.

\section{References}\label{references}

{[}OAI{]} OpenAI, \emph{Uniform exclusion of Landau--Siegel zeros}, manuscript dated 1 October 2026. \href{https://github.com/openai/math/blob/adc7f1241b42e322a6451854ab7e4b4c146bf78a/preprints/Uniform-exclusion-of-Landau-Siegel-zeros-October-1-2026/build/paper.tex}{Pinned manuscript source}; \href{https://github.com/openai/math/blob/adc7f1241b42e322a6451854ab7e4b4c146bf78a/preprints/Uniform-exclusion-of-Landau-Siegel-zeros-October-1-2026/paper.pdf}{manuscript PDF record}. Used: Theorem 1 and Sections 2--6, particularly the evaluation determinant and weighted row selection. TeX source consulted directly.

{[}Ba{]} William D. Banks, \emph{On the OpenAI proof that Landau--Siegel zeros do not exist}, arXiv:2610.15483v1, 11 October 2026. \href{https://arxiv.org/html/2610.15483v1}{Text}. Used: introduction, Theorem 1.1, proof organization and authorship disclosures. Text consulted directly.

{[}Ke{]} Bryce Kerr, \emph{A simpler proof that there are no Landau--Siegel zeros}, arXiv:2610.21099v1, 15 October 2026. \href{https://arxiv.org/html/2610.21099v1}{Text}. Used: the pigeonhole/norm construction and Lemmas 5--7, including the statement attributed to Masser. Text consulted directly.

{[}Gr{]} Ben Green, \emph{On Landau--Siegel zeros}, hexagon:2610.00311v2, October 2026. \href{https://hexagonmath.org/2610.00311v2}{Record}; \href{https://hexagonmath.org/pdf/2610.00311v2}{text}. Used: Theorem 1.1, Sections 2--6 and Appendix A. Version-2 text and provenance record consulted directly.

{[}Br{]} Kevin Barreto, \emph{An elementary proof of the non-existence of Landau--Siegel zeros}, three-page manuscript supplied as page images. \textbf{Document supplied by @AcerFur, via the \href{https://x.com/AcerFur/status/2112239827077026128}{X post identified by the submitter}.} Used: Theorem 1 and its proof on pp.~1--2; Corollary 2 and its proof on pp.~2--3; opening attribution and bibliography. The images were read directly. The sharing URL was supplied by the submitter; the post's contents were not independently retrieved. No separate repository identifier, publication date or licence for the manuscript is inferred. This submission summarizes the argument and does not reproduce the page images.

{[}LZZ{]} Rick F. Lu, Asif Zaman and Haonan Zhao, \emph{Numerical Computations concerning Landau--Siegel Zeros}, arXiv:2602.03626v1, 3 February 2026. \href{https://arxiv.org/html/2602.03626v1}{Text}. Used: Theorem 1.1, the finite-conductor result invoked by Banks. Statement consulted; computation not rerun.

{[}Ma{]} D. W. Masser, \emph{On polynomials and exponential polynomials in several complex variables}, Inventiones Mathematicae 63 (1981), 81--95. \href{https://link.springer.com/article/10.1007/BF01389194}{Publisher record}. Bibliographic record checked; the result used in the comparison was consulted through Kerr's Lemma 6, not independently reread in the original paper.

{[}MV{]} Hugh L. Montgomery and Robert C. Vaughan, \emph{Multiplicative Number Theory I: Classical Theory}, Cambridge Studies in Advanced Mathematics 97, Cambridge University Press, 2006. \href{https://www-personal.umich.edu/~hlm/mnt1.html}{Author's book page}. Bibliographic record checked. Barreto cites Corollary 2.6, Theorem 2.7(b), Theorem 9.13 and Corollary 10.18 for the classical inputs; these numbered references are reported from his manuscript, not claimed to have been independently rechecked in the book.

{[}WS{]} We-Solve, \href{https://we-solve.github.io/problems/siegel-zeros/}{Math-77: Siegel Zeros} and \href{https://we-solve.github.io/contributions/siegel-zeros/attempt-2/r1/}{Math-77-A2: OpenAI claim summary}. Used: exact problem scope, existing literature contribution and its verification caveat. Pages consulted directly on 19 October 2026.

\end{document}
